\section{Estadística descriptiva}
\label{sec:descriptiva}

La estadística descriptiva complementa el análisis exploratorio mediante
medidas de posición, dispersión y forma. Mientras la
Sección~\ref{sec:eda} identificó patrones generales de la cartera, esta
sección cuantifica esos patrones y examina con mayor detalle la
distribución del monto del crédito.

Para las variables numéricas se calcularon la media, la desviación
estándar muestral, la mediana, los cuartiles, el rango intercuartílico,
los valores mínimo y máximo, la asimetría y la curtosis en exceso. La
media y la desviación estándar permiten resumir el centro y la dispersión
cuando la distribución no presenta una asimetría extrema; la mediana y
el rango intercuartílico ofrecen una descripción más robusta ante colas
largas y observaciones alejadas del centro.

\subsection{Resumen de las variables numéricas}

La Tabla~\ref{tab:descriptivos-globales} presenta los principales
estadísticos de las cinco variables numéricas originales.

\begin{table}[H]
    \centering
    \caption{Estadística descriptiva de las variables numéricas.}
    \label{tab:descriptivos-globales}

    \footnotesize
    \resizebox{\textwidth}{!}{
        \begin{tabular}{
            @{}
            l
            r
            r
            r
            r
            r
            r
            r
            r
            @{}
        }
            \toprule
            Variable
            & Media
            & Desv.\ est.
            & Mediana
            & Q1
            & Q3
            & RIC
            & Mín.--Máx.
            & Asimetría \\
            \midrule

            Duración
            & \num{20.90}
            & \num{12.06}
            & \num{18}
            & \num{12}
            & \num{24}
            & \num{12}
            & \num{4}--\num{72}
            & \num{1.09} \\

            Monto
            & \num{3271.26}
            & \num{2822.74}
            & \num{2319.50}
            & \num{1365.50}
            & \num{3972.25}
            & \num{2606.75}
            & \num{250}--\num{18424}
            & \num{1.95} \\

            Edad
            & \num{35.55}
            & \num{11.38}
            & \num{33}
            & \num{27}
            & \num{42}
            & \num{15}
            & \num{19}--\num{75}
            & \num{1.02} \\

            Créditos en el banco
            & \num{1.41}
            & \num{0.58}
            & \num{1}
            & \num{1}
            & \num{2}
            & \num{1}
            & \num{1}--\num{4}
            & \num{1.27} \\

            Dependientes
            & \num{1.16}
            & \num{0.36}
            & \num{1}
            & \num{1}
            & \num{1}
            & \num{0}
            & \num{1}--\num{2}
            & \num{1.91} \\

            \bottomrule
        \end{tabular}
    }

    \begin{minipage}{\textwidth}
        \footnotesize
        \textit{Nota.} \(Q1\) y \(Q3\) representan los cuartiles
        primero y tercero. RIC corresponde al rango intercuartílico.
        La duración se expresa en meses, el monto en marcos alemanes
        (DM) y la edad en años. Todas las variables tienen
        \(n=1000\) observaciones y no presentan valores faltantes.
    \end{minipage}
\end{table}

La duración media es de \num{20.90} meses, mientras que su mediana es de
\num{18} meses. La diferencia entre ambas medidas y una asimetría de
\num{1.09} muestran que los plazos se concentran en valores relativamente
bajos, pero incluyen una cola formada por créditos de mayor duración.

La edad presenta un comportamiento similar. Su media es de
\num{35.55} años y su mediana de \num{33} años, con una asimetría
positiva de \num{1.02}. Aunque la mayor parte de los solicitantes se
concentra entre \num{27} y \num{42} años, la muestra incluye personas de
hasta \num{75} años.

Las variables \variable{n_creditos_banco} y
\variable{n_dependientes} tienen una escala reducida. En particular, el
rango intercuartílico de \variable{n_dependientes} es cero porque los
cuartiles primero, segundo y tercero toman el valor uno. Esto no significa
ausencia total de variabilidad, sino una fuerte concentración en la
categoría más frecuente.

\subsection{Forma de la distribución del monto}

El monto es la variable con mayor heterogeneidad. La media es de
\num{3271.26} DM, mientras que la mediana es de \num{2319.50} DM. La
diferencia de casi mil marcos entre ambas medidas evidencia la influencia
de los créditos situados en la cola derecha.

La variable presenta:

\[
\operatorname{asimetría}=\num{1.95},
\qquad
\operatorname{curtosis\ en\ exceso}=\num{4.29}.
\]

La asimetría positiva indica que la distribución se extiende hacia los
montos elevados. La curtosis en exceso sugiere una concentración central
y unas colas más pronunciadas que las de una distribución normal. El
monto máximo, \num{18424} DM, es casi ocho veces la mediana de la
muestra.

Estos resultados justifican tres decisiones posteriores:

\begin{enumerate}
    \item complementar la media con la mediana y el rango
    intercuartílico;

    \item utilizar pruebas robustas o no paramétricas al comparar
    grupos;

    \item considerar el logaritmo del monto cuando se requiera reducir
    la asimetría o modelar relaciones multiplicativas.
\end{enumerate}

La transformación logarítmica no elimina la información de los créditos
de mayor monto. Su función es comprimir la escala y reducir la influencia
relativa de las observaciones extremas.

\subsection{Descriptivos por clase de riesgo}

La Tabla~\ref{tab:descriptivos-clase} compara duración, monto y edad
entre los grupos de buen y mal riesgo.

\begin{table}[H]
    \centering
    \caption{Estadística descriptiva de las variables principales según
    la clase de riesgo.}
    \label{tab:descriptivos-clase}

    \footnotesize
    \resizebox{\textwidth}{!}{
        \begin{tabular}{
            @{}
            l
            l
            r
            r
            r
            r
            r
            r
            r
            @{}
        }
            \toprule
            Variable
            & Clase
            & \(n\)
            & Media
            & Desv.\ est.
            & Q1
            & Mediana
            & Q3
            & RIC \\
            \midrule

            Duración
            & Buen riesgo
            & 700
            & \num{19.21}
            & \num{11.08}
            & \num{12}
            & \num{18}
            & \num{24}
            & \num{12} \\

            Duración
            & Mal riesgo
            & 300
            & \num{24.86}
            & \num{13.28}
            & \num{12}
            & \num{24}
            & \num{36}
            & \num{24} \\

            \addlinespace

            Monto
            & Buen riesgo
            & 700
            & \num{2985.46}
            & \num{2401.47}
            & \num{1375.50}
            & \num{2244}
            & \num{3634.75}
            & \num{2259.25} \\

            Monto
            & Mal riesgo
            & 300
            & \num{3938.13}
            & \num{3535.82}
            & \num{1352.50}
            & \num{2574.50}
            & \num{5141.50}
            & \num{3789} \\

            \addlinespace

            Edad
            & Buen riesgo
            & 700
            & \num{36.22}
            & \num{11.38}
            & \num{27}
            & \num{34}
            & \num{42.25}
            & \num{15.25} \\

            Edad
            & Mal riesgo
            & 300
            & \num{33.96}
            & \num{11.22}
            & \num{25}
            & \num{31}
            & \num{40}
            & \num{15} \\

            \bottomrule
        \end{tabular}
    }

    \begin{minipage}{\textwidth}
        \footnotesize
        \textit{Nota.} Los resultados son descriptivos y no constituyen
        por sí solos evidencia de diferencias poblacionales. Los
        contrastes formales se presentan en la
        Sección~\ref{sec:inferencia}.
    \end{minipage}
\end{table}

La diferencia más visible corresponde a la duración. Los malos créditos
tienen una mediana de \num{24} meses, frente a \num{18} meses en los
buenos. Además, su rango intercuartílico es el doble, lo que muestra una
mayor variabilidad de los plazos dentro del grupo de mal riesgo.

El monto medio aumenta de \num{2985.46} DM a \num{3938.13} DM y la
mediana de \num{2244} DM a \num{2574.50} DM. La diferencia entre medias
es mayor que la diferencia entre medianas, lo que confirma que parte del
contraste está influido por operaciones elevadas en la cola derecha.

La edad muestra una diferencia más moderada. Los solicitantes de mal
riesgo tienen una mediana de \num{31} años, frente a \num{34} años en
los buenos créditos. El solapamiento entre grupos continúa siendo amplio,
por lo que ninguna de estas variables debe interpretarse de manera
aislada como una regla de clasificación.

\subsection{Ajuste paramétrico del monto mediante máxima verosimilitud}

Debido a que el monto es positivo y presenta una cola derecha
pronunciada, se compararon dos familias continuas con soporte positivo:
Lognormal y Gamma. Ambas distribuciones se ajustaron mediante máxima
verosimilitud utilizando las funciones
\texttt{scipy.stats.lognorm.fit()} y
\texttt{scipy.stats.gamma.fit()}.

En los dos modelos se fijó el parámetro de localización en cero:

\[
\texttt{loc}=0.
\]

Esta restricción mantiene el soporte de las distribuciones a partir de
cero y deja dos parámetros libres en cada modelo.

Para la distribución Lognormal:

\[
X \sim \operatorname{Lognormal}(\mu,\sigma),
\]

de modo que:

\[
\log(X)\sim\mathcal{N}(\mu,\sigma^2).
\]

Con \(\texttt{loc}=0\), los estimadores cerrados son:

\[
\widehat{\mu}
=
\frac{1}{n}
\sum_{i=1}^{n}\log(x_i),
\]

\[
\widehat{\sigma}
=
\sqrt{
    \frac{1}{n}
    \sum_{i=1}^{n}
    \left[
        \log(x_i)-\widehat{\mu}
    \right]^2
}.
\]

La estimación obtenida fue:

\[
\widehat{\mu}=\num{7.788691},
\qquad
\widehat{\sigma}=\num{0.776086}.
\]

Los valores coinciden, hasta la precisión numérica, con los parámetros
devueltos por SciPy. Esta comprobación sirve como verificación
independiente de la implementación.

Para la distribución Gamma se utilizó la parametrización forma--escala:

\[
X\sim\operatorname{Gamma}(k,\theta),
\]

con estimaciones:

\[
\widehat{k}=\num{1.7921},
\qquad
\widehat{\theta}=\num{1825.3962}.
\]

\subsection{Comparación mediante log-verosimilitud y AIC}

Los modelos se compararon mediante la log-verosimilitud maximizada y el
criterio de información de Akaike:

\[
\operatorname{AIC}
=
2k-2\ell(\widehat{\theta}),
\]

donde \(k\) es el número de parámetros libres y
\(\ell(\widehat{\theta})\) la log-verosimilitud evaluada en los
estimadores de máxima verosimilitud. En ambos modelos se utilizó
\(k=2\), porque el parámetro de localización se mantuvo fijo.

\begin{table}[H]
    \centering
    \caption{Comparación de los ajustes Lognormal y Gamma para el monto
    del crédito.}
    \label{tab:comparacion-mle}

    \begin{tabular}{
        @{}
        l
        l
        r
        l
        r
        r
        r
        @{}
    }
        \toprule
        Distribución
        & Parámetro 1
        & Valor
        & Parámetro 2
        & Valor
        & Log-verosimilitud
        & AIC \\
        \midrule

        Lognormal
        & \(\mu\)
        & \num{7.7887}
        & \(\sigma\)
        & \num{0.7761}
        & \num{-8954.1375}
        & \num{17912.2751} \\

        Gamma
        & \(k\)
        & \num{1.7921}
        & \(\theta\)
        & \num{1825.3962}
        & \num{-9007.2152}
        & \num{18018.4305} \\

        \bottomrule
    \end{tabular}
\end{table}

La distribución Lognormal presenta una log-verosimilitud mayor y un AIC
menor. La diferencia es:

\[
\Delta\operatorname{AIC}
=
\num{18018.4305}
-
\num{17912.2751}
=
\num{106.1554}.
\]

Dentro de las dos familias evaluadas, esta diferencia favorece de manera
clara al modelo Lognormal. El resultado no demuestra que la Lognormal sea
la distribución verdadera del monto ni que proporcione un ajuste
perfecto. El AIC permite realizar una comparación relativa entre los
modelos candidatos, pero no constituye una prueba absoluta de bondad de
ajuste.

La evidencia gráfica complementaria muestra que los cuantiles de
\(\log(\variable{monto})\) siguen aproximadamente una línea recta en la
parte central, aunque presentan desviaciones en las colas. En
consecuencia, la Lognormal ofrece una aproximación útil para describir la
forma general de la distribución, pero conserva limitaciones en los
valores extremos.

\begin{hallazgo}
Entre las distribuciones comparadas, la Lognormal describe mejor el
monto del crédito. Este resultado respalda el uso de
\variable{log_monto} para reducir la asimetría, pero no sustituye los
diagnósticos específicos de los modelos posteriores.
\end{hallazgo}

\subsection{Implicancias para los análisis posteriores}

La transformación logarítmica se empleará cuando resulte conveniente
para reducir la escala de los montos elevados y representar relaciones
proporcionales. En el modelo OLS, su utilidad deberá evaluarse mediante
el comportamiento de los residuos, la homocedasticidad y la validación
cruzada.

En la regresión logística, la transformación no se utiliza porque los
predictores deban seguir una distribución normal. El modelo logístico no
exige normalidad marginal de las variables explicativas. Su incorporación
responde a una decisión sobre la forma funcional de la relación entre el
monto y el logit de la probabilidad de mal crédito.

Finalmente, la asimetría observada justifica que la comparación del monto
entre buenos y malos créditos no se limite a una prueba sobre medias. En
la Sección~\ref{sec:inferencia} se complementará la prueba de Welch con
un contraste de Mann--Whitney basado en rangos.

% Contenido previsto:
% - Resumen global y por clase.
% - Media, mediana, desviación, cuartiles, IQR, asimetría y curtosis.
% - Ajuste MLE lognormal y gamma del monto.
% - Comparación mediante log-verosimilitud y AIC.
